\appendix

\section{超参数}

与 GLM-5 模型架构相关的超参数见 \Cref{tab:model_arch_glm45_glm5}。

在训练方面，我们沿用 GLM-4.5 的设置，包括 Muon 优化器、余弦衰减和 batch size warmup。学习率先经历从 0 到 2e-4 的 warmup 阶段，再衰减至 4e-5 直至预训练结束。在中期训练阶段，学习率从 4e-5 线性下降到 1e-5。其余超参数与 GLM-4.5 相同。对于 DSA warmup 阶段，学习率从 5e-3 下降到 2e-4。对于 DSA 稀疏适配阶段，我们使用恒定学习率 1e-5。

\begin{table}[!ht]
    \centering
    \caption{GLM-4.5 与 GLM-5 的模型架构。参数统计中，所有模型均计入 MTP 层参数，但不计入词嵌入与输出层参数。}
    \label{tab:model_arch_glm45_glm5}
    \begin{tabular}{l|cc}
    \toprule
       Model & \textbf{GLM-4.5} & \textbf{GLM-5} \\
    \midrule
       \# Total Parameters & 355B & 744B \\
       \# Activated Parameters & 32B & 40B \\
       \# Dense Layers & 3 & 3 \\
       \# MoE Layers & 89 & 75 \\
       \# MTP Layers & 1 & 1 \\
       Hidden Dim & 5120 & 6144 \\
       Dense Intermediate Dim & 12288 & 12288 \\
       MoE Intermediate Dim & 1536 & 2048 \\
    \midrule
       QK Head Dim & 128 & 192 \\
       V Head Dim & 128 & 256 \\
       Q LoRA Dim & -- & 2048 \\
       KV LoRA Dim & -- & 512 \\
       \# Attention Heads & 96 & 64 \\
       \# Key-Value Heads & 8 & -- \\
       \# Indexer Attn Heads & -- & 32 \\
       \# Indexer Head Dim & -- & 128 \\
    \midrule
       \# Experts (total) & 160 & 256 \\
       \# Routed Experts & 8 & 8 \\
       \# Shared Experts & 1 & 1 \\
    \midrule
       Vocabulary Size & 151552 & 154880 \\
    \bottomrule
    \end{tabular}
\end{table}

\section{评测细节}
\subsection{基座模型评测}
我们在 \Cref{tab:base_benchmark} 中使用英文、中文、代码和数学基准评测 GLM-5 基座模型。

\begin{table}[!ht]
  \centering
  \caption{GLM-5-Base、GLM-4.5-Base 与其他代表性开源基座模型对比。}
  \resizebox{\textwidth}{!}{
  \begin{tabular}{lcccc|c}
    \toprule
     & \textbf{Benchmark (Metric)} & \textbf{DeepSeek-V3} & \textbf{Kimi-K2} & \textbf{GLM-4.5} & \textbf{GLM-5} \\
    & & \textbf{Base} & \textbf{Base} & \textbf{Base} & \textbf{Base} \\
    \midrule
    \multirow{3}{*}{} & Architecture & MoE & MoE & MoE & MoE \\
    & \# Activated Params & 37B & 32B & 32B & 40B \\
    & \# Total Params & 671B & 1043B & 355B & 744B \\
    \midrule
    \multirow{6}{*}{\textbf{English}} & SimpleQA (EM) & 26.6 & 35.3 & 30.0 & 36.0 \\
    & BBH (EM) & 88.4 & 88.7 & 86.2 & 87.4 \\
    & MMLU (EM) & 87.2 & 87.8 & 86.1 & 88.3 \\
    & HellaSwag (EM) & 88.9 & 94.6 & 87.1 & 88.1 \\
    & PIQA (EM) & 84.7 & - & 85.3 & 84.6 \\
    & TriviaQA (EM) & 82.9 & 85.1 & 80.0 & 80.9 \\
    \midrule
    \multirow{2}{*}{\textbf{Code}} & EvalPlus (Pass@1) & 65.6 & 80.3 & 78.1 & 87.0 \\
    & LiveCodeBench-Base (Pass@1) & 24.6 & 26.3 & 28.1 & 34.4 \\
    \midrule
    \multirow{2}{*}{\textbf{Math}} & GSM8K (EM) & 87.6 & 92.1 & 79.4 & 68.8 \\
    & MATH (EM) & 62.6 & 70.2 & 61.0 & 56.4 \\
    \midrule
    \multirow{4}{*}{\textbf{Chinese}} & CLUEWSC (EM) & 82.7 & - & 83.5 & 84.2 \\
    & C-Eval (EM) & 90.1 & 92.5 & 86.9 & 88.8 \\
    & C3 (EM) & 78.6 & - & 83.1 & 80.3 \\
    & Chinese-SimpleQA (EM) & 72.1 & 77.6 & 70.1 & 74.6 \\
    \bottomrule
  \end{tabular}
  }
  \label{tab:base_benchmark}
\end{table}

\subsection{ARC 基准评测}
\label{sec:eval_details_arc}
Humanity’s Last Exam（HLE）与其他推理任务：我们设置最大生成长度为 $131,072$ tokens（$temperature=1.0, top\_p=0.95, max\_new\_tokens=131072$）。默认报告文本子集；标记 * 的结果来自全量集合。我们使用 GPT-5.2（medium）作为 judge 模型。对 HLE-with-tools，最大上下文长度设为 $202,752$ tokens。

SWE-bench \& SWE-bench Multilingual：我们在 OpenHands 中使用定制指令提示运行 SWE-bench 套件。设置为：$temperature=0.7, top\_p=0.95, max\_new\_tokens=16384$，上下文窗口 200K。

BrowseComp：不使用上下文管理时，保留最近 5 轮细节。使用上下文管理时，采用与 DeepSeek-V3.2 和 Kimi K2.5 相同的 discard-all 策略。

Terminal-Bench 2.0（Terminus 2）：我们使用 Terminus 框架评测，设置为 $timeout=2h, temperature=0.7, top\_p=1.0, max\_new\_tokens=8192$，上下文窗口 128K。资源上限为 16 CPUs 与 32 GB RAM。

Terminal-Bench 2.0（Claude Code）：我们在 Claude Code 2.1.14（think mode）中评测，设置为 $temperature=1.0, top\_p=0.95, max\_new\_tokens=65536$。我们移除墙钟时限，同时保留单任务 CPU 与内存约束。我们修复了 Claude Code 引入的环境问题，并额外报告修复歧义指令的 verified Terminal-Bench 2.0 数据集结果（见：\url{https://huggingface.co/datasets/zai-org/terminal-bench-2-verified}）。分数取 5 次运行均值。

CyberGym：我们在 Claude Code 2.1.18（think mode，无 web 工具）中评测，设置为（$temperature=1.0, top\_p=1.0, max\_new\_tokens=32000$），每任务超时 250 分钟。结果为 1,507 个任务的单次 Pass@1。

MCP-Atlas：全部模型均在 think mode 下于 500 任务公开子集上评测，单任务超时 10 分钟。我们使用 Gemini 3 Pro 作为 judge 模型。

$\tau^2$-Bench：我们在 Retail 和 Telecom 中加入小幅提示调整，以避免过早用户终止导致失败。对于 Airline，采用 Claude Opus 4.5 system card 提出的领域修复。

Vending-Bench 2：运行由 Andon Labs 独立完成\footnote{\url{https://andonlabs.com/evals/vending-bench-2}}。


\subsection{$\tau^2$-Bench 的优化用户模拟器}
\label{sec:tau_user_prompt}
我们在 Telecom 与 Retail 中加入小幅提示调整，以避免过早用户终止导致失败。优化后的提示见图~\ref{fig:tau2-prompt-telecom} 与图~\ref{fig:tau2-prompt-retail}。这些优化提示被如下集成到系统提示中：

\begin{tcolorbox}[left=0mm,right=0mm,top=0mm,bottom=0mm,boxsep=1mm,arc=0mm,boxrule=0pt, frame empty, breakable]
\begin{lstlisting}
SYSTEM_PROMPT = """"
{global_user_sim_guidelines}

<scenario>
{instructions}
</scenario>

{optimized_user_prompt}
"""".strip()
\end{lstlisting}
\end{tcolorbox}

\begin{tcolorbox}[left=0mm,right=0mm,top=0mm,bottom=0mm,boxsep=1mm,arc=0mm,boxrule=0pt, frame empty, breakable]
\footnotesize
\begin{lstlisting}
# Note:
- Do not generate the '###TRANSFER###' before agent clearly tells "YOU ARE BEING TRANSFERRED TO A HUMAN AGENT. PLEASE HOLD ON.".
    Example:
    Case1:
        - agent: "Would you like me to transfer you to a human agent who can assist you with these options and help get your service restored?"
        - user: "Yes, please transfer me to a human agent.".
    Case2:
        - user: "YOU ARE BEING TRANSFERRED TO A HUMAN AGENT. PLEASE HOLD ON."
        - user: "###TRANSFER###"
\end{lstlisting}
\end{tcolorbox}
\noindent\begin{minipage}{\textwidth}
\captionof{figure}{$\tau^2$-Bench Telecom 的优化用户提示。}
\label{fig:tau2-prompt-telecom}
\end{minipage}

\begin{tcolorbox}[left=0mm,right=0mm,top=0mm,bottom=0mm,boxsep=1mm,arc=0mm,boxrule=0pt, frame empty, breakable]
\footnotesize
\begin{lstlisting}
# Rules:
- Just generate one line at a time to simulate the user's message.
- Do not give away all the instruction at once. Only provide the information that is necessary for the current step.
- Do not hallucinate information that is not provided in the instruction. Follow these guidelines:
    1. If the agent asks for information NOT in the instruction:
        - Say you don't remember or don't have it
        - Offer alternative information that IS mentioned in the instruction
    2. Examples:
        - If asked for order ID (not in instruction): "Sorry, I don't remember the order ID, can you search for it? My name/email/phone number/zipcode is ..."
        - If asked for email (not in instruction): "I don't have my email handy, but I can give you my name and zip code which are..."
- Do not repeat the exact instruction in the conversation. Instead, use your own words to convey the same information.
- Try to make the conversation as natural as possible, and stick to the personalities in the instruction.
# Constraint Handling:
- Provide requests strictly based on what is explicitly stated in the instruction.
- Do not assume, extend, substitute, or generalize in any form.
- Do not modify or relax constraints on:
- Time / Date
- Budget
- Specific terms (e.g., "same" must not be replaced with "similar")
- Core Rule: Any attribute NOT mentioned in the instruction can be either changed or kept the same
- Examples:
    - If instruction says "exchange red item to blue": Only color must change, other attributes (size, material, etc.) are flexible
    - If instruction says "exchange red item to blue, keep the same size": Both color must change AND size must stay the same
- Exception: Only follow additional constraints when explicitly stated in the instruction
# Domain-Specific Rules:
## For Retail scenarios:
- Focus on product attributes and exchange/return processes as specified in instructions.
- During confirmations: Always respond based strictly on the original instruction, never deviate to match agent's provided options. Restate your requirement from the instruction rather than selecting from agent's choices.
    - Example: If the agent provides specific options (A/B/C) but the instruction states a general requirement (e.g., "same as pending order"), always restate or confirm based on what the instruction says, not by directly selecting from the agent's provided options.
# When NOT to finish the conversation:
- Do not end until you have clearly and completely expressed all your requirements and constraints.
- Do not end until the agent has completed all tasks mentioned in the instruction and verified no operations were missed.
- Do not end if the agent's execution results do not match your expectations or are incorrect/incomplete.
# When you CAN finish the conversation:
- Only when all above conditions are satisfied AND all tasks are completed correctly.
- OR when you have clearly expressed complete requirements but the system explicitly states it cannot complete them due to technical limitations - in this case, accept transfer to human.
# How to finish the conversation:
- If the agent has completed all tasks, generate the '###STOP###' token to end the conversation.
# Note:
- You should carefully check if the agent has completed all tasks mentioned in the instruction before generating '###STOP###'.
\end{lstlisting}
\end{tcolorbox}
\noindent\begin{minipage}{\textwidth}
\captionof{figure}{$\tau^2$-Bench Retail 的优化用户提示。}
\label{fig:tau2-prompt-retail}
\end{minipage}

\subsection{真实智能体工程体验评估}

\subsubsection{前端评估}

\label{sec:frontend_eval_data_cons}

\paragraph{数据。}
我们的数据集包含 7 类前端场景，用于评估模型在不同功能域下的工程能力：Business Management Systems、Web Games、SVG/Canvas Rendering、Creative Tools \& Editors、Showcase Pages、Forms \& Tables、Data Visualization。
% \begin{itemize}
% \item Business Management Systems: Applications for enterprise or personal data/workflow management.
% \item Web Games: Interactive, real-time rendering-focused entertainment applications.
% \item SVG/Canvas Rendering: Graphics-intensive applications based on low-level drawing APIs.
% \item Creative Tools \& Editors: Utility applications supporting content creation and complex state management.
% \item Showcase Pages: Layout-oriented pages emphasizing visual expression and responsiveness.
% \item Forms \& Tables: Applications centered on structured data entry, validation, and processing.
% \item Data Visualization: Graphical representation and analytical tools for complex datasets.
% \end{itemize}


\begin{table}[h]
\centering
\small
\caption{前端应用场景分布。}
\begin{tabular}{lp{7cm}rr}
\toprule
\textbf{Category} & \textbf{Description} & \textbf{\# Tasks} & \textbf{\# Checkitems} \\
\midrule
Business Systems & 企业/个人数据与流程管理。 & 42 & 167 \\
Web Games & 以交互与娱乐为核心的游戏应用。 & 40 & 163\\
SVG/Canvas & 图形渲染与交互可视化。 & 32 & 166\\
Creative Tools & 内容创作与在线编辑工具。 & 28 & 160\\
Showcase Pages & 视觉表达与信息展示页面。 & 27 & 115\\
Forms \& Tables & 结构化数据录入与处理。 & 26 & 93\\
Data Visualization & 图形化数据表达与分析。 & 25 & 85\\
\bottomrule
\end{tabular}

\end{table}

\subparagraph{按编程技术栈的数据分布}
该基准完整覆盖三类主流范式：Vanilla Web Stack（HTML/CSS/JS）、React 组件化框架、Vue 3 + Vite 渐进式方案。
\begin{table}[h]
\centering
\caption{技术栈与评估单元统计。}
\begin{tabular}{llcc}
\toprule
\textbf{Category} & \textbf{Description} & \textbf{\# Tasks} & \textbf{\# Checkitems} \\
\midrule
HTML & 原生 HTML/CSS/JS 开发。 & 113 & 490 \\
React & 组件化框架开发。 & 58 & 249 \\
Vue & Vue 3 + Vite 渐进式方案。 & 49 & 210 \\
\bottomrule
\end{tabular}
\end{table}

\subparagraph{数据样例}
每个测试样例由三部分组成：\textbf{Task}、\textbf{Checklist} 与 \textbf{专用环境}。下方是一个代表性样例：
\begin{tcolorbox}[left=0mm,right=0mm,top=0mm,bottom=0mm,boxsep=1mm,arc=0mm,boxrule=0pt, frame empty, breakable]
\footnotesize
\begin{lstlisting}
    Task: Develop an online drawing tool that includes a brush, an eraser, a white canvas, and a save button.
        The brush color and thickness should be selectable via buttons on the left. Users can draw on the canvas by clicking and dragging the mouse.
        The eraser size should be selectable via buttons on the left. Users can erase content by clicking and dragging the mouse over the canvas.
        Once the drawing is complete, clicking the "Save" button should allow the user to save the image locally.
        Please implement this using the React framework in the current directory.

    Checklist:
        The user can select the brush color and thickness using the left-hand buttons, and drawing is functional via mouse click-and-drag on the canvas.
        The user can select the eraser size using the left-hand buttons, and erasing is functional via mouse click-and-drag on the canvas.
        Upon clicking the "Save" button, the generated image is successfully saved to the local machine.
\end{lstlisting}
\end{tcolorbox}
\noindent\begin{minipage}{\textwidth}
\end{minipage}


\subparagraph{\textbf{数据构建与验证}}

我们实现了严格的四阶段流程以确保数据质量：
\begin{itemize}[leftmargin=1.5em,itemsep=2pt]
\item 阶段 1：任务合成。由资深前端专家设计任务，确保贴近真实工程挑战，同时在场景与技术上保持分布均衡。
\item 阶段 2：检查项生成与精炼。我们先使用 Claude Sonnet 4.5 基于任务规格 $T$ 合成候选 checklist，再由专家细致审核与整合。经多轮打磨，确保每个 check-item 语义无歧义、客观可判，并对用户需求形成充分覆盖。
\item 阶段 3：基于执行的纠偏。我们在 Agent-as-a-Judge 框架与人工专家之间进行交叉验证。任何判定分歧都会触发底层数据复评与修正，以消除潜在噪声。
\item 阶段 4：动态基准迭代。为维持区分度，我们迭代更新测试集，移除已无法挑战 SOTA 编码智能体的简单任务。该专家主导整理流程最终形成 220 个高质量前端编码任务及其对应 checklist。
\end{itemize}
